%% ============================================================
%%  Signal N3: The Conditional DVOL Signal
%%  Compile: pdflatex A1.tex  (twice for refs)
%% ===================================================================
\documentclass[12pt, a4paper]{article}

\usepackage{amsmath, amssymb, amsthm, mathtools}
\usepackage{booktabs}
\usepackage{geometry}
\usepackage{hyperref}
\usepackage{xcolor}
\usepackage{array}
\usepackage{multirow}
\usepackage{enumitem}
\usepackage{caption}
\usepackage{float}
\usepackage{parskip}
\usepackage{bm}

\geometry{margin=2.5cm}

\hypersetup{colorlinks=true, linkcolor=blue!60!black, urlcolor=blue!60!black}

\theoremstyle{plain}
\newtheorem{theorem}{Theorem}[section]
\newtheorem{proposition}{Proposition}[section]
\newtheorem{corollary}{Corollary}[section]
\theoremstyle{definition}
\newtheorem{definition}{Definition}[section]
\newtheorem{remark}{Remark}[section]

%% Macros
\newcommand{\IC}{\mathrm{IC}}
\newcommand{\E}{\mathbb{E}}
\newcommand{\Var}{\mathrm{Var}}
\newcommand{\Cov}{\mathrm{Cov}}
\newcommand{\Cor}{\mathrm{Cor}}
\newcommand{\rank}{\mathrm{rank}}
\newcommand{\DVOL}{\mathrm{DVOL}}
\newcommand{\iid}{\overset{\text{iid}}{\sim}}
\newcommand{\bp}{\,\mathrm{bp}}

\title{\textbf{Signal N3: The Conditional DVOL Signal}\\[0.5em]
\large Fear-Resolution Return Predictability in BTCUSDT Perpetual Futures\\[0.3em]
\large A Complete Treatment from First Principles}
\author{}
\date{May 2026}

\begin{document}
\maketitle
\tableofcontents
\newpage

%% ============================================================
\section{Introduction}
%% ============================================================

This report documents the design, derivation, and rigorous empirical validation
of Signal N3: a return-predictive signal for the BTCUSDT USDT-margined
perpetual futures contract on Binance, constructed from the Deribit implied
volatility index (DVOL).

The core result is this: the 30-day $z$-score of Deribit DVOL
predicts the next 24-hour log-return of BTCUSDT with a Pearson
Information Coefficient of $+0.133$ on a non-overlapping daily sample
of 365 out-of-sample days ($p = 0.017$, block-bootstrap with 14-day blocks).
The $|\IC|/\IC^*$ ratio of $4.87\times$ the maker-cost breakeven means
the signal generates, in expectation, nearly five times the edge needed to cover
round-trip transaction costs. A threshold-filtered long rule produces Sharpe $= 2.95$ over 863 OOS days
(2024--2026-YTD) with maximum drawdown $-1{,}612\bp$.
\emph{(Entry thresholds are withheld from this public version.  All
statistical results are reported in full.)}

The report proceeds as follows. Section~\ref{sec:market} introduces the
relevant market structures from first principles: perpetual futures mechanics
and the Deribit DVOL index. Section~\ref{sec:signal} defines the signal.
Section~\ref{sec:stats} derives the statistical machinery from scratch:
log-returns, the Pearson IC, the breakeven IC proposition (with full proof),
the overlapping-return inflation problem, and the block bootstrap.
Section~\ref{sec:results} presents all empirical results on the 2024
out-of-sample period. Section~\ref{sec:kill} is the 2025--2026 kill-attempt.
Section~\ref{sec:regime} develops the conditional strategy. Section~\ref{sec:backtest}
presents the position-level backtest. Section~\ref{sec:discussion} discusses
mechanisms, limitations, and multiple-comparison concerns.


%% ============================================================
\section{Market Background}
\label{sec:market}
%% ============================================================

\subsection{Perpetual Futures}

A \emph{perpetual futures contract} on an asset $S$ is a derivative that
trades like a futures contract but has no expiry date. The contract tracks
a spot index price $I_t$ (the weighted average of BTC/USDT prices across
major spot exchanges) through a mechanism called the \emph{funding rate}.

\paragraph{The funding rate.}
Every 8 hours (at 00:00, 08:00, and 16:00 UTC), a cash payment changes hands
between all long and all short position holders. If the perpetual mark price
$M_t$ has been trading above the spot index $I_t$ --- meaning longs are
overenthusiastic relative to spot --- then longs pay shorts:
\[
  F_t = \underbrace{\frac{M_t - I_t}{I_t}}_{\text{premium}}
        + \text{clamp}\!\left(r_{\text{interest}} - \frac{M_t - I_t}{I_t},\,
                              -0.05\%,\, 0.05\%\right)
\]
where $r_{\text{interest}}$ is a fixed per-settlement interest differential
(Binance uses $r_{\text{interest}} = 0.01\%/8\text{h}$, i.e.\ $0.03\%$ per
calendar day across three daily settlements, for BTCUSDT).
When $F_t > 0$, a long position of notional $N$ pays $N \times F_t$ per
settlement. When $F_t < 0$, shorts pay longs.

The funding mechanism is a continuous arbitrage force: if the perpetual
persistently trades above spot, long holders are steadily charged,
incentivising them to close longs (buying pressure dissipates) or short
sellers to sell the perpetual (arbitrageurs short the perp and buy spot
until prices converge). The funding rate therefore proxies the
\emph{aggregate leverage imbalance} in the market.

\paragraph{Total PnL for a perpetual long.}
For a long position of unit notional opened at price $C_\tau$ and closed
at $C_{\tau+h}$, the total P\&L over the hold period is:
\[
  \pi_\tau^{(h)} = \underbrace{\log\frac{C_{\tau+h}}{C_\tau}}_{r_\tau^{(h)}}
                  - \underbrace{\sum_{k:\, \tau_k \in [\tau, \tau+h]} F_{\tau_k}}_{f_\tau^{(h)}}
\]
The first term is the price log-return; the second is the sum of three
funding payments within the 24-hour window (at positive funding, the long
pays; at negative funding, the long receives). The signal N3 is evaluated
on total PnL $\pi_\tau^{(h)}$ as well as on price return alone
(Section~\ref{sec:totalPnL}).

\subsection{The Deribit DVOL Index}
\label{sec:dvol}

\subsubsection{Model-Free Implied Variance}

For any traded underlying with options, the implied volatility of a single
option at strike $K$ and maturity $T$ is the unique $\sigma$ such that the
Black--Scholes formula reprices the observed option mid-quote. However,
individual implied volatilities depend on the model assumed and vary
across the strike surface (the ``volatility smile''), making them
difficult to aggregate.

The \emph{model-free implied variance}, introduced by
Britten-Jones \& Neuberger (2000) and refined for discrete strikes
by Jiang \& Tian (2005), is a strike-aggregated measure of the
risk-neutral expectation of realised variance over horizon $T$:

\begin{definition}[Model-Free Implied Variance]
Under no-arbitrage and continuity of the underlying, the
\emph{annualised} risk-neutral expected variance --- the risk-neutral
expectation of integrated variance divided by the horizon $T$ --- can be
recovered from option prices without specifying a model for the underlying:
\[
  \frac{1}{T}\,\E^Q_t\!\left[\int_t^{t+T} \sigma_s^2\,\mathrm{d}s\right]
  = \frac{2\,\mathrm{e}^{rT}}{T}\!\left[
    \int_0^F \frac{P_t(K)}{K^2}\,\mathrm{d}K
    + \int_F^\infty \frac{C_t(K)}{K^2}\,\mathrm{d}K
  \right]
\]
where $F$ is the at-the-money forward price, $C_t(K)$ is the time-$t$
price of a call with strike $K$, $P_t(K)$ is the put price, and
$\mathrm{e}^{rT}$ accounts for discounting.  Dividing by $T$ converts
the total integrated variance (units: variance $\times$ time) into an
annualised rate directly comparable across maturities and consistent
with the discrete approximation in \eqref{eq:vix}.
\end{definition}

In practice, options are quoted at a finite grid of strikes
$\{K_1 < K_2 < \cdots < K_M\}$, and the integrals are replaced by
a trapezoidal sum (CBOE, 2009):
\begin{equation}
  \hat{\sigma}^2_{T} = \frac{2\,\mathrm{e}^{rT}}{T}\,
                       \sum_{i=1}^{M} \frac{\Delta K_i}{K_i^2}\,Q(K_i)
                       - \frac{1}{T}\!\left(\frac{F}{K_0} - 1\right)^{\!2}
  \label{eq:vix}
\end{equation}
where $Q(K_i)$ is the mid-quote of the out-of-the-money option at strike
$K_i$, $\Delta K_i = (K_{i+1} - K_{i-1})/2$ is the half-spacing, and
$K_0$ is the at-the-money strike.

\subsubsection{Deribit BTCVOL (DVOL)}

Deribit computes BTCVOL (DVOL) using~\eqref{eq:vix} applied to BTC
options expiring approximately 30 calendar days from now, interpolating
between two adjacent maturities to achieve an exact 30-day target.
The index is reported in annualised percentage units:
\[
  \DVOL_t = 100 \times \sqrt{\hat{\sigma}^2_{30\text{d}} \times 365}
\]
A DVOL of 50 means that the options market prices the 30-day annualised
BTC volatility at 50\%, implying a 1-day standard deviation of
$50\% / \sqrt{365} \approx 2.6\%$ and a 30-day standard deviation of
$50\%/\sqrt{365/30} \approx 14.4\%$.

DVOL is published at hourly frequency. Over the sample period (2022--2026),
DVOL ranged from approximately 30 to 110, with a mean of 52 and standard
deviation of 12 (approximately, based on the distribution used in
Table~\ref{tab:dvol_dist}).

The key property exploited in this paper: \textbf{DVOL is a
market-consensus summary of near-term uncertainty}, aggregated across all
strike prices and option tenors near 30 days, without imposing a model for
the underlying BTC price process.


%% ============================================================
\section{Signal Construction}
\label{sec:signal}
%% ============================================================

\subsection{Three DVOL Variants}

Let $\DVOL_\tau$ denote the hourly DVOL observation at hour $\tau$.
Three signals are constructed from the same source:

\begin{definition}[N1: DVOL Level]
  $N1_\tau = \DVOL_\tau$, the raw implied volatility index.
\end{definition}

\begin{definition}[N2: 8-Hour Slope]
  $N2_\tau = \DVOL_\tau - \DVOL_{\tau - 8}$, the change in DVOL
  over the preceding 8 hours.
\end{definition}

\begin{definition}[N3: 30-Day $z$-Score]
  \[
    N3_\tau = \frac{\DVOL_\tau - \mu_{30\mathrm{d},\tau}}
                   {\sigma_{30\mathrm{d},\tau}}
  \]
  where $\mu_{30\mathrm{d},\tau} = \frac{1}{720}\sum_{k=0}^{719}\DVOL_{\tau-k}$
  and $\sigma_{30\mathrm{d},\tau}$ is the corresponding sample standard
  deviation over the rolling 720-hour (30-day) window.
\end{definition}

All three signals are then forward-filled to the 1-minute bar grid
($\DVOL$ is updated at hourly frequency, so all 60 minute-bars within
an hour share the same signal value) and merged with the BTCUSDT 1-minute
kline dataset.

\subsection{Why the $z$-Score Dominates N1 and N2}
\label{sec:whyz}

N1 captures the \emph{level} of implied volatility. BTC perpetuals have a
structural tendency to have higher realised volatility in bull markets than
bear markets (amplified call demand), so a high N1 level may merely reflect
a high-volatility market regime rather than an elevated-fear state relative
to that regime.

N2 captures the \emph{direction of change} in IV over 8 hours. This proxies
the momentum of fear, but it ignores the starting level and confounds the
sign of IV changes (a fall from 80 to 72 and a fall from 40 to 32 produce
the same N2 despite very different regime states).

N3 corrects for both of these shortcomings by expressing DVOL as a
\emph{deviation from its own recent history}, standardised by recent
variability. Two properties follow immediately:

\begin{enumerate}[noitemsep]
  \item \textbf{Regime-relative interpretation}: N3 = 1.5 means DVOL is
        1.5 standard deviations above its 30-day mean, regardless of whether
        the absolute level is 45 or 75.
  \item \textbf{Stationarity}: unlike N1, which has a clear long-term upward
        trend with BTC maturation and exhibits non-stationary level changes
        across years, N3 is approximately mean-zero and unit-variance at all
        times (by construction), making IC estimates comparable across periods.
\end{enumerate}

The statistical consequence (Section~\ref{sec:incremental}): after
orthogonalising N3 against N1 (regressing out the level), the residual IC
is $+0.133$ ($4.87\times$ maker breakeven) --- identical to the unconditional
N3 IC. N1 carries essentially no incremental predictive information beyond
what N3 already encodes.

\subsection{The Fear-Resolution Hypothesis}
\label{sec:mechanism}

The sign of the N3 IC is \emph{positive}: a high $z$-score (DVOL above its
30-day mean) predicts positive 24-hour returns. This is the central finding
and requires a mechanism.

\paragraph{Hypothesis: Fear peaks are contrarian predictors.}
When DVOL spikes above its local mean, the options market is pricing elevated
near-term uncertainty. This elevated IV reflects (and causes) the following
market dynamics:

\begin{enumerate}
  \item \textbf{Options dealers are short gamma.} Options dealers who have
        sold near-the-money options are short gamma: their delta hedge requires
        them to sell as BTC falls and buy as BTC rises. This amplifies intraday
        price moves during fear events.

  \item \textbf{Leveraged longs liquidate.} Elevated volatility triggers
        margin calls and forces liquidation of levered positions. Perpetual
        funding rates often fall (or invert) as longs are flushed.

  \item \textbf{Retail shorts pile in.} Elevated fear attracts momentum
        short sellers who expect further decline.

  \item \textbf{Fear resolves stochastically.} The fear state does not persist
        indefinitely. As the triggering event passes or uncertainty resolves,
        the options market reprices, DVOL mean-reverts toward its baseline, and
        the three effects above unwind simultaneously: dealers
        reduce their hedges (less selling pressure), new longs replace the
        liquidated ones, and short sellers cover.
\end{enumerate}

The \emph{resolution} of a fear peak therefore generates positive price
pressure. Critically, this mechanism is only available to trade when DVOL is
elevated above its local mean ($N3 > 0$) --- it is absent in calm environments
($N3 \approx 0$).

\paragraph{Empirical predictions of the hypothesis.}
If the fear-resolution mechanism is the correct explanation, then:

\begin{enumerate}[label=(\alph*)]
  \item IC $\approx 0$ in low-DVOL regimes (no fear to resolve). \quad \textbf{Confirmed: $p = 0.691$ for DVOL $< 51$.}
  \item IC strong in high-DVOL regimes. \quad \textbf{Confirmed: $p = 0.018$ for DVOL $\geq 51$.}
  \item IC predicts direction more than magnitude. \quad \textbf{Confirmed: $4.87\times$ vs.\ $2.03\times$.}
  \item IC robust to removing extreme event days. \quad \textbf{Confirmed: $4.09\times$ ex-event days.}
  \item Training IC below OOS peak (causal not spurious). \quad \textbf{Confirmed: Train $2.11\times$ vs.\ OOS-H1 $7.05\times$.}
\end{enumerate}

All five predictions are borne out by the data, providing strong circumstantial
evidence that the mechanism is structural rather than a spurious correlation.


%% ============================================================
\section{Statistical Framework}
\label{sec:stats}
%% ============================================================

\subsection{Log-Returns as Prediction Targets}

For a price series $\{C_t\}$, the \emph{simple return} over a horizon $h$ is
$\tilde{r}_t^{(h)} = C_{t+h}/C_t - 1$. The \emph{log-return} is:
\[
  r_t^{(h)} = \log\frac{C_{t+h}}{C_t} = \log C_{t+h} - \log C_t
\]
Log-returns are preferred for three reasons:

\paragraph{Temporal additivity.}
Log-returns over sub-periods add exactly:
$r_{t,t+h+k} = r_{t,t+h} + r_{t+h,t+h+k}$.
Simple returns compound multiplicatively, requiring chain multiplication.
This additivity is essential for equity curve construction
(the log-equity is the cumulative sum of log-returns).

\paragraph{Symmetry.}
A simple return of $+50\%$ is not cancelled by a subsequent $-50\%$:
$1.5 \times 0.5 = 0.75 \neq 1$. Log-returns are symmetric: $+\log 1.5$
is cancelled by $-\log 1.5$. This symmetry makes long and short positions
analytically equivalent, since a short log-return is $-r_t^{(h)}$.

\paragraph{Cost comparability.}
Transaction costs are approximately constant fractions of notional regardless
of direction. A round-trip cost $c_{\text{RT}}$ is subtracted uniformly from
$r_t^{(h)}$ regardless of its sign or magnitude. This is exact in log-return
space and approximate in simple return space.

Throughout this report, $r_t^{(h)} \in \mathbb{R}$ denotes the 1-minute-bar
forward log-return at horizon $h \in \{1\text{h}, 4\text{h}, 8\text{h},
24\text{h}, 48\text{h}\}$.

\paragraph{Basis-point convention.}
All PnL, drawdown, and per-trade figures are stated in \emph{basis points}
using the universal financial convention: $1\bp = 0.01\% = 10^{-4}$.
Concretely, a log-return of $0.01$ (1\%) equals $100\bp$, and a
cumulative log-return of $0.93$ (93\%) equals $9{,}300\bp$.
This is the decimal return multiplied by $10^{4}$; the code enforces
this via \texttt{net\_ret * 1e4}.
The abbreviation ``bp'' never means $0.1\%$ in this document.

\subsection{The Information Coefficient}
\label{sec:IC}

\begin{definition}[Pearson IC]
  Let $\{(s_t, r_t)\}_{t=1}^n$ be a sample of signal-return pairs.
  The \emph{Pearson Information Coefficient} is the linear correlation
  between the signal and the return:
  \begin{equation}
    \IC(s, r) = \frac{\displaystyle\sum_{t=1}^n (s_t - \bar{s})(r_t - \bar{r})}
                     {\displaystyle\sqrt{\sum_{t=1}^n (s_t - \bar{s})^2}\;
                      \sqrt{\sum_{t=1}^n (r_t - \bar{r})^2}}
    \label{eq:pearson_ic}
  \end{equation}
  where $\bar{s}$ and $\bar{r}$ are the sample means.
\end{definition}

\paragraph{Why Pearson?}
The IC* breakeven Proposition~\ref{prop:ICstar} is derived under the
linear signal model
\[
  r_t^{(h)} = \IC \cdot \sigma_r \cdot z_t + \varepsilon_t,
\]
in which $\IC$ is by definition the Pearson correlation between $z_t$ and
$r_t^{(h)}$.  Consistency between the empirical IC measure and the
theoretical breakeven formula requires Pearson throughout: applying
the breakeven formula to a Spearman IC would compare quantities derived
from different definitions of correlation.

\begin{remark}[Spearman as robustness check]
  The \emph{Spearman IC} applies the Pearson formula to the rank vectors
  $R_t = \rank(s_t)$, $Q_t = \rank(r_t)$:
  \[
    \IC_S(s, r) = \frac{\sum_t (R_t - \bar{R})(Q_t - \bar{Q})}
                       {\sqrt{\sum_t (R_t - \bar{R})^2}
                        \sqrt{\sum_t (Q_t - \bar{Q})^2}}
    = 1 - \frac{6\sum_t d_t^2}{n(n^2-1)}, \quad d_t = R_t - Q_t
  \]
  It is invariant to monotone transformations and more robust to outliers.
  For bivariate Gaussian $(s,r)$,
  $\IC_S \approx (6/\pi)\arcsin(\IC_P/2)$,
  which implies $|\IC_S - \IC_P| \approx 4$--$5\%$ for the IC magnitudes
  observed here (e.g.\ $\IC_P = 0.133 \Rightarrow \IC_S \approx 0.127$).
  Spearman ICs are reported in robustness tables for comparison.
\end{remark}

\begin{remark}
  A negative IC is equally informative as a positive IC of the same magnitude:
  the signal should simply be traded on the opposite side. We report
  $|\IC|$ and the sign separately throughout.
\end{remark}

\subsection{The Breakeven Information Coefficient}
\label{sec:ICstar}

A statistically significant IC does not imply profitability. A signal must
overcome round-trip transaction costs. We now derive the minimum IC required.

\begin{proposition}[Breakeven IC]
  \label{prop:ICstar}
  Consider a trading strategy that enters long when a standardised signal
  $z_t \sim \mathcal{N}(0,1)$ is positive and short when $z_t$ is negative,
  exiting after $h$ bars.  Under the linear-IC model
  \begin{equation}
    r_t^{(h)} = \IC \cdot \sigma_r \cdot z_t + \varepsilon_t, \qquad
    \varepsilon_t \perp z_t,\quad
    \Var(r_t^{(h)}) = \sigma_r^2
    \label{eq:linearIC}
  \end{equation}
  the expected gross profit per trade equals
  $\IC \cdot \sigma_r \cdot \sqrt{2/\pi}$.  The signal is profitable
  against a round-trip cost $c_{\mathrm{RT}}$ if and only if
  \begin{equation}
    |\IC| > \IC^* \;\coloneqq\; \frac{c_{\mathrm{RT}}}{\sigma_r \sqrt{2/\pi}}
    \label{eq:ICstar}
  \end{equation}
\end{proposition}

\begin{proof}
  Under model~\eqref{eq:linearIC}, since $z_t \perp \varepsilon_t$:
  \[
    \E[r_t^{(h)} \mid z_t] = \IC \cdot \sigma_r \cdot z_t
  \]
  For a long trade conditional on $z_t > 0$:
  \begin{align*}
    \E[r_t^{(h)} \mid z_t > 0]
    &= \IC \cdot \sigma_r \cdot \E[z_t \mid z_t > 0]
    = \IC \cdot \sigma_r \cdot \frac{\E[z_t \cdot \mathbf{1}_{z_t > 0}]}
                                     {P(z_t > 0)}
  \end{align*}
  For $z_t \sim \mathcal{N}(0,1)$:
  \[
    \E[z_t \mid z_t > 0]
    = \frac{\int_0^\infty z\,\phi(z)\,\mathrm{d}z}{\tfrac{1}{2}}
    = 2\int_0^\infty z\,\phi(z)\,\mathrm{d}z
    = 2\left[-\phi(z)\right]_0^\infty
    = 2\phi(0)
    = \sqrt{\tfrac{2}{\pi}}
  \]
  By symmetry, a short trade ($z_t < 0$) earns $\E[-r_t^{(h)} \mid z_t < 0]
  = \IC \cdot \sigma_r \cdot \sqrt{2/\pi}$ as well.

  Expected net profit per trade equals gross minus round-trip cost:
  \[
    \E[\text{net profit}] = \IC \cdot \sigma_r \sqrt{2/\pi} - c_{\mathrm{RT}}
  \]
  Setting this to zero and solving for IC gives~\eqref{eq:ICstar}.\qed
\end{proof}

\begin{remark}[Threshold entry]
  If the strategy only trades when $|z_t| > c$ for some $c > 0$, the expected
  gross profit per trade is $\IC \cdot \sigma_r \cdot \E[z \mid z > c]$. Since
  $\E[z \mid z > c] = \phi(c)/[1-\Phi(c)] > \sqrt{2/\pi}$ for all $c > 0$,
  the breakeven IC is \emph{lower} than~\eqref{eq:ICstar}. The formula
  as stated is conservative: it gives the threshold for the most
  disadvantageous trading rule (enter on every bar).
\end{remark}

\begin{remark}[Consistency with the empirical IC]
  Proposition~\ref{prop:ICstar} is derived for the Pearson IC, which is
  also the empirical measure used throughout (Section~\ref{sec:IC}).
  The two are therefore consistent: the breakeven threshold $\IC^*$ and
  the observed $\IC$ are both Pearson correlations.

  For reference, the Spearman approximation $\IC_S \approx (6/\pi)\arcsin(\IC_P/2)$
  gives $\IC_S \approx \IC_P \times 0.955$ for the IC magnitudes in this report
  (a 4--5\% downward bias relative to Pearson).  Using Spearman ICs with the
  same $\IC^*$ formula would therefore understate the $|\IC|/\IC^*$ ratios
  by approximately 5\%.
\end{remark}

Table~\ref{tab:ICstar} gives the breakeven IC for each prediction horizon,
using maker-side costs ($c_{\text{RT}}^{\text{maker}} = 0.06\%$ round-trip).

\begin{table}[H]
\centering
\caption{Breakeven IC thresholds at each horizon (maker cost $c_{\mathrm{RT}} = 0.06\%$)}
\label{tab:ICstar}
\begin{tabular}{lrr}
\toprule
Horizon $h$ & $\sigma_r^{(h)}$ & $\IC^*_{\mathrm{maker}}$ \\
\midrule
1h   & $0.561\%$ & $0.134$ \\
4h   & $1.104\%$ & $0.068$ \\
8h   & $1.549\%$ & $0.049$ \\
24h  & $2.680\%$ & $0.028$ \\
48h  & $3.720\%$ & $0.020$ \\
\bottomrule
\end{tabular}
\end{table}

The longer the horizon, the lower $\IC^*$, because a larger $\sigma_r$ means
the strategy earns more gross profit per trade from the same IC. This is why
24-hour and 48-hour signals are much easier to trade profitably than
1-hour signals, even if their ICs are similar in magnitude.

\subsection{The Overlapping-Return Inflation Problem}
\label{sec:overlap}

The 1-minute BTCUSDT dataset has $T \approx 527{,}040$ observations in the OOS
2024 window. If we naively compute $r_t^{(24\text{h})} = \log(C_{t+1440}/C_t)$
for every $t$ and treat the $527{,}040$ pairs $(N3_t, r_t^{(24\text{h})})$ as
independent observations, we will severely overstate statistical significance.

\paragraph{The autocorrelation structure of overlapping returns.}
Under the simplifying assumption that 1-minute log-returns are i.i.d.\ with
variance $\delta^2$, the 24h log-return decomposes as:
\[
  r_t^{(24\text{h})} = \sum_{j=0}^{1439} \Delta_j(t),
  \qquad \Delta_j(t) = \log\frac{C_{t+j+1}}{C_{t+j}}
\]
Two adjacent observations share all but one minute of their window:
\[
  r_{t+1}^{(24\text{h})} - r_t^{(24\text{h})} = \Delta_{1439}(t) - \Delta_0(t)
\]
which is the net effect of adding one new minute and removing one old minute.
More generally, for lag $k \leq 1440$:
\begin{align*}
  \Cov\!\left(r_t^{(24\text{h})}, r_{t+k}^{(24\text{h})}\right)
  &= \Cov\!\left(\sum_{j=0}^{1439}\Delta_j(t),\,
                 \sum_{j=0}^{1439}\Delta_j(t+k)\right)
   = (1440 - k)\,\delta^2
\end{align*}
since the two sums share exactly $1440 - k$ common terms (for $k < 1440$).
The autocorrelation at lag $k$ is therefore:
\begin{equation}
  \Cor\!\left(r_t^{(24\text{h})}, r_{t+k}^{(24\text{h})}\right)
  = \frac{(1440 - k) \delta^2}{1440\,\delta^2}
  = 1 - \frac{k}{1440}, \qquad 0 \leq k \leq 1440
  \label{eq:overlap_autocor}
\end{equation}
This is a tent function that linearly decays from 1 at lag 0 to 0 at lag 1440.

\paragraph{Effective sample size.}
For a stationary series with autocorrelation function $\rho_k$, the effective
sample size for variance estimation is approximately:
\[
  n_{\mathrm{eff}} \approx \frac{T}{1 + 2\sum_{k=1}^{\infty} \rho_k}
\]
Substituting~\eqref{eq:overlap_autocor}:
\[
  1 + 2\sum_{k=1}^{1440-1}\!\!\left(1 - \frac{k}{1440}\right)
  = 1 + 2 \cdot \frac{1440 \cdot 1439}{2 \cdot 1440}
  = 1 + 1439
  = 1440
\]
Therefore:
\begin{equation}
  n_{\mathrm{eff}} \approx \frac{T}{1440} = \frac{527{,}040}{1440} \approx 366
  \label{eq:neff}
\end{equation}
The effective number of independent 24-hour return observations is
approximately 366 days --- one per day --- regardless of whether we
compute the return at every 1-minute bar or sample once per day.

\paragraph{Practical implication.}
The $t$-statistic for the observed IC from the full $n = 527{,}040$ series is
inflated by a factor of $\sqrt{T / n_{\mathrm{eff}}} = \sqrt{1440} \approx 38$
relative to the correct statistic based on $n_{\mathrm{eff}} = 366$. An IC
that appears to have $p \approx 10^{-20}$ using the naive formula has a true
$p$ that could be anywhere from $0.001$ to $0.50$ depending on the signal's
autocorrelation structure.

\textbf{The solution adopted in this paper}: sample every 1440 bars
($n = 365$ non-overlapping daily observations per OOS year) and apply the
block bootstrap to control for any residual autocorrelation.

\subsection{The Block Bootstrap}
\label{sec:bootstrap}

The block bootstrap (Künsch, 1989; Liu \& Singh, 1992) provides valid inference
for statistics of dependent time series without parametric assumptions about
the dependence structure.

\begin{definition}[Non-Overlapping Block Bootstrap]
  Given a time series $(X_1, \ldots, X_n)$ and a statistic $T_n$:
  \begin{enumerate}
    \item Choose block length $\ell$.  Form $m = \lfloor n/\ell \rfloor$
          complete non-overlapping blocks
          $B_j = (X_{(j-1)\ell+1}, \ldots, X_{j\ell})$ for $j = 1, \ldots, m$.
          The tail $r = n - m\ell < \ell$ observations are discarded; for the
          primary parameters ($n = 365$, $\ell = 14$) this is $r = 1$ day
          (0.27\% of the sample), which has negligible effect on inference.
    \item Draw $\lceil n/\ell \rceil$ blocks
          $B_{j_1}^*, \ldots, B_{j_{\lceil n/\ell\rceil}}^*$ uniformly
          at random \emph{with replacement} from $\{B_1, \ldots, B_m\}$.
          The same block may appear more than once; this is intentional ---
          repeated blocks generate the variability in the bootstrap null
          distribution.
    \item Concatenate the drawn blocks, then trim to exactly $n$ observations
          to form the bootstrap resample $X^*$.
    \item Compute the bootstrap statistic $T_n^* = T_n(X^*)$.
    \item Repeat $B$ times. The empirical distribution of
          $\{T_n^{*(b)}\}_{b=1}^B$ approximates the sampling distribution
          of $T_n$.
  \end{enumerate}
\end{definition}

For inference on $\IC$, the block bootstrap preserves temporal autocorrelation
within blocks (capturing weekly seasonality, funding-cycle autocorrelation,
and momentum effects) while permuting the dependence structure across blocks.

\paragraph{Block length selection.}
The block length $\ell$ should exceed the effective autocorrelation decay
length of the signal-return interaction. For the daily non-overlapping series
used here, a block of $\ell = 14$ days (two calendar weeks) captures the
relevant dependence: weekly trading patterns, bi-weekly options gamma pinning
cycles near options expiry, and the 8-day funding reset cycle.

\paragraph{$p$-value computation.}
For a two-sided test of $H_0 : \IC = 0$, we compute the permutation $p$-value:
\[
  p = \frac{1}{B}\sum_{b=1}^B \mathbf{1}\!\left[|T_n^{*(b)}| \geq |\widehat{\IC}|\right]
\]
In all bootstrap tests reported, $B = 1{,}000$ replicates are used.

\subsection{Walk-Forward Out-of-Sample Validation}
\label{sec:walkforward}

Standard cross-validation violates temporal causality: using future data to
estimate a model evaluated on the past. For time-series prediction, the
correct protocol is \emph{walk-forward} validation:

\begin{enumerate}[noitemsep]
  \item Fix all signal parameters using only the training period
        (2023-01-01 to 2023-12-31).
  \item Evaluate the signal \emph{without any re-fitting} on disjoint
        OOS windows: 2024-H1, 2024-H2, 2025-H1, 2025-H2, 2026-YTD.
  \item A signal that shows near-zero IC on training data but strong IC
        on OOS data cannot be the result of overfitting the training data
        (it would have shown up there first).
\end{enumerate}

For Signal N3, no parameters are fit on the training data: the 30-day rolling
window and $z$-score transformation are fixed by definition; the period
DVOL regime filter threshold is the only tuned parameter, and it
is cross-validated in Section~\ref{sec:threshold_sweep}.


%% ============================================================
\section{Empirical Results: OOS 2024}
\label{sec:results}
%% ============================================================

All results in this section use the out-of-sample period
2024-01-01 to 2024-12-31 ($n_{\text{daily}} = 365$ non-overlapping
daily observations) unless stated otherwise. Signal N3 was constructed
using only data up to the evaluation date; no OOS data leaks into the
signal construction.

\subsection{Overlap Inflation Check}

Before presenting the main results, we verify that the non-overlapping
sampling does not \emph{reduce} the observed IC relative to the full
overlapping series (i.e., that we are not throwing away information ---
we are just removing spurious significance).

\begin{table}[H]
\centering
\caption{Overlap inflation sanity check (OOS 2024, N3 $\to r^{(24h)}$)}
\label{tab:overlap}
\begin{tabular}{lrrl}
\toprule
Series & $n$ & IC (Pearson) & Ratio ($|\IC|/\IC^*$) \\
\midrule
1m-overlapping (all bars)      & 525,600 & $+0.134$ & $4.78\times$ \\
Non-overlapping (daily sample) & 365     & $+0.133$ & $4.87\times$ \\
\midrule
Inflation factor (overlap / non-overlap) & \multicolumn{3}{l}{$0.98\times$ (no inflation)} \\
\bottomrule
\end{tabular}
\end{table}

The non-overlapping series gives a \emph{higher} IC than the overlapping
series. This means the non-overlapping test is not conservative: the true
predictive relationship is fully preserved (indeed, marginally strengthened)
when we restrict to one observation per day. There is no evidence of
overlap inflation in this signal.

Why? Equation~\eqref{eq:overlap_autocor} showed that adjacent daily returns
share 1439/1440 of their window. However, the N3 signal is a 30-day
$z$-score updated hourly, so adjacent 1-minute N3 values within the same
hour are \emph{identical}. The within-hour signal repetition introduces a
compensating downward bias to the overlapping IC. The two effects roughly
cancel, leaving no net inflation.

\subsection{Block Bootstrap on Daily Non-Overlapping Series}
\label{sec:boot_results}

\begin{table}[H]
\centering
\caption{Block-bootstrap significance test: non-overlapping daily observations
($B=1000$, block $=14$ days, OOS 2024). IC computed via Pearson correlation
(Section~\ref{sec:IC}); Spearman values are $\approx\!4$--$5\%$ lower.}
\label{tab:n3_rigorous}
\begin{tabular}{llrrrrl}
\toprule
Signal & Hz & IC (Pearson) & $p$ & Null 95\% range$^\dagger$ & $|\IC|/\IC^*$ & Verdict \\
\midrule
N3 ($z$-score) & 24h & $+0.133$ & $\mathbf{0.017}$ & $[-0.099, +0.113]$ & $4.87\times$ & \textbf{Pass} \\
N3 ($z$-score) & 48h & $+0.164$ & $\mathbf{0.002}$ & $[-0.102, +0.105]$ & $8.29\times$ & \textbf{Pass} \\
N2 (slope)     & 24h & $+0.100$ & $0.055$           & $[-0.099, +0.105]$ & $3.67\times$ & Fail ($p$) \\
N2 (slope)     & 48h & $+0.099$ & $0.050$           & $[-0.099, +0.097]$ & $4.98\times$ & Fail ($p$) \\
N1 (level)     & 24h & $+0.019$ & $0.698$           & $[-0.098, +0.109]$ & $0.70\times$ & Fail \\
N1 (level)     & 48h & $+0.004$ & $0.931$           & $[-0.095, +0.109]$ & $0.21\times$ & Fail \\
\bottomrule
\end{tabular}
\smallskip

\begin{minipage}{0.92\textwidth}
\footnotesize $^\dagger$ \textbf{This is not a confidence interval for IC.}
The column reports the 2.5th--97.5th percentiles of the \emph{bootstrap
null distribution}: $B = 1000$ ICs obtained by permuting the signal
within temporal blocks while holding returns fixed.  Under the null
of no predictability, a randomly generated IC falls in this range
95\% of the time.  The p-value is the fraction of null ICs with
$|\mathrm{IC}_\text{null}| \geq |\mathrm{IC}_\text{obs}|$.
For N3 at 24h: $\mathrm{IC}_\text{obs} = +0.133$ lies above the
null upper bound $+0.113$, consistent with $p = 0.017$.
\end{minipage}
\end{table}

\textbf{N3 is the only DVOL variant that passes rigorous daily testing.}
N1 and N2 both fail the block-bootstrap $p < 0.05$ criterion at all horizons.
N3 passes at both 24h ($p = 0.017$) and 48h ($p = 0.002$). The 48h result
at $8.29\times$ breakeven and $p = 0.002$ is especially strong; it is
evaluated as a secondary confirmation rather than the primary trading horizon
since holding 48 hours doubles the funding exposure and is practically
harder to implement cleanly.

\subsection{Signal Correlations and Incremental IC}
\label{sec:incremental}

The Spearman correlation matrix among N1, N2, N3 (computed on daily
non-overlapping OOS):
\[
  \begin{pmatrix}
    \rho(\text{N1},\text{N1}) & \rho(\text{N1},\text{N2}) & \rho(\text{N1},\text{N3}) \\
    \rho(\text{N2},\text{N1}) & \rho(\text{N2},\text{N2}) & \rho(\text{N2},\text{N3}) \\
    \rho(\text{N3},\text{N1}) & \rho(\text{N3},\text{N2}) & \rho(\text{N3},\text{N3})
  \end{pmatrix}
  =
  \begin{pmatrix}
    1.00 & 0.07 & 0.56 \\
    0.07 & 1.00 & 0.21 \\
    0.56 & 0.21 & 1.00
  \end{pmatrix}
\]

N3 and N1 are moderately correlated ($\rho = 0.56$): a high $z$-score
tends to occur when the absolute level is also elevated. N2 and N1
are nearly uncorrelated ($\rho = 0.07$).

\paragraph{Incremental IC via orthogonalisation.}
To test whether N3 carries predictive information \emph{beyond} N1, we
project N3 onto N1 via ordinary least squares and take the residual:
\[
  \tilde{N3} = N3 - \hat{\beta}_0 - \hat{\beta}_1 \cdot N1
\]
The IC of this residual $\tilde{N3}$ against $r^{(24h)}$ measures the
information in N3 not already captured by N1.

Result: IC($\tilde{N3}$, $r^{(24h)}$) $= +0.133$ ($4.87\times$).
Removing N1's contribution does not reduce N3's IC.
\textbf{N1 adds nothing to N3; N3 is the only signal required.}

\subsection{Walk-Forward IC by Half-Year}
\label{sec:walkforward_ic}

\begin{table}[H]
\centering
\caption{N3 walk-forward IC by period (non-overlapping daily, 24h horizon).
  The $|\IC_P|/\IC^*$ ratio uses each period's own realised $\sigma_r^{(24h)}$
  in~\eqref{eq:ICstar}, not the pooled value from Table~\ref{tab:ICstar}.}
\label{tab:wf_2024}
\begin{tabular}{lrrrr}
\toprule
Period & $n$ & IC (Pearson) & IC (Spearman) & $|\IC_P|/\IC^*$ \\
\midrule
Train 2023   & 365 & $+0.069$ & $+0.005$ & $2.11\times$ \\
OOS 2024-H1  & 182 & $+0.188$ & $+0.177$ & $7.05\times$ \\
OOS 2024-H2  & 183 & $+0.075$ & $+0.091$ & $2.68\times$ \\
\bottomrule
\end{tabular}
\end{table}

The training Pearson IC of $2.11\times$ is positive but substantially lower
than OOS 2024-H1 ($7.05\times$), and the signal has no parameters fit
to the training data (N3 is a fixed z-score formula with no tuned thresholds
at this stage). The positive training IC reflects genuine structural
predictability, not overfitting. Notably, the Spearman training IC is near
zero ($+0.005$) while the Pearson IC is $+0.069$: the linear relationship
is driven by a small number of extreme N3 observations that correctly
predicted large positive returns, consistent with the threshold-rule design
(only extreme fear events are traded).

\subsection{Total PnL Including Funding}
\label{sec:totalPnL}

For a long position in the perpetual, the total return includes the funding
payment received or paid at each settlement. To evaluate whether the N3
edge is eroded by funding costs, we define:
\[
  \pi_t^{(24h)} = r_t^{(24h)} - f_t^{(24h)}
\]
where $f_t^{(24h)} = \sum_{k: \tau_k \in [t, t+1440]} F_{\tau_k}$ is the
total funding paid by a long over the 24-hour horizon (approximately three
settlements). For 2024, the mean funding rate was negative (shorts paying
longs on average), at $-3.3\bp/\text{day}$.

\begin{table}[H]
\centering
\caption{N3 IC: price return vs.\ total PnL including funding (OOS 2024)}
\label{tab:funding}
\begin{tabular}{lrr}
\toprule
Target & IC (Pearson) & Ratio \\
\midrule
Price return only $r^{(24h)}$       & $+0.133$ & $4.87\times$ \\
Total PnL $\pi^{(24h)}$ (price + funding) & $+0.129$ & $4.74\times$ \\
\midrule
Mean daily funding cost (long, 2024) & \multicolumn{2}{l}{$-3.3\bp/\text{day}$ (negative: shorts paid longs)} \\
\bottomrule
\end{tabular}
\end{table}

The IC is essentially unchanged by funding inclusion. In 2024, the funding
rate was on average negative (shorts paying longs), so the perpetual long
received a positive carry. The N3 edge is robust to the full P\&L measure.

\subsection{Net-of-Cost Quintile Returns}

Sort OOS 2024 daily observations into quintiles by N3 signal value.
Compute the mean total 24h PnL (price + funding) within each quintile,
then subtract maker round-trip cost.

\begin{table}[H]
\centering
\caption{N3 quintile total PnL at 24h horizon, net of maker cost
(OOS 2024, $n=365$ days)}
\label{tab:quintile}
\begin{tabular}{lrrrr}
\toprule
Quintile & $n$ & Mean N3$z$ & Gross 24h (bp) & Net (maker) (bp) \\
\midrule
Q1 (lowest N3)  & 73 & $-1.76$ & $-17.7$ & $-23.7$ \\
Q2              & 73 & $-0.84$ & $-12.8$ & $-18.8$ \\
Q3              & 73 & $-0.07$ & $+1.9$  & $-4.1$  \\
Q4              & 73 & $+0.78$ & $+41.7$ & $+35.7$ \\
Q5 (highest N3) & 73 & $+2.00$ & $+78.0$ & $+72.0$ \\
\midrule
Q5$-$Q1 spread  &    &         & $+95.7$ & $+83.7$ \\
\bottomrule
\end{tabular}
\end{table}

The quintile structure is monotone and economically large. Q5 earns
$+72.0\bp/\text{day}$ net after maker costs. Q1 earns $-23.7\bp/\text{day}$
(a profitable short). The Q5--Q1 net spread of $83.7\bp/\text{day}$ represents
the full long-short edge.

Three features of Table~\ref{tab:quintile} are noteworthy:
\begin{enumerate}[noitemsep]
  \item The relationship is monotone across all five quintiles, suggesting
        the predictive relationship is not driven by a single tail.
  \item The spread between Q4 and Q3 ($39.8\bp$) is larger than Q3--Q2
        ($14.7\bp$), which is larger than Q2--Q1 ($5.1\bp$): the effect
        is convex in the signal, concentrated in the upper quintiles.
  \item Q5 net PnL exceeds Q1 net PnL in absolute value ($72.0 > 23.7$),
        indicating the long side of the trade is stronger than the short side.
\end{enumerate}

\subsection{Long-Short Rule Performance}
\label{sec:longshort}

Define a parameterised long-short rule with threshold $\theta \geq 0$:
\[
  \text{pos}_t = \begin{cases}
    +1 & \text{if } N3_t > \theta \\
    -1 & \text{if } N3_t < -\theta \\
    \hspace{2.5mm}0 & \text{otherwise}
  \end{cases}
\]
Daily rebalancing, maker round-trip cost charged on position changes.

\begin{table}[H]
\centering
\caption{Long-short rule, varying $\theta$ (OOS 2024, $n=365$ daily, total PnL, maker cost)}
\label{tab:longshort}
\begin{tabular}{rrrrrrr}
\toprule
$\theta$ & Active days & Trades & Net PnL (bp) & Sharpe & MaxDD (bp) \\
\midrule
0.00 & 365 & 42  & $+8{,}591$  & $1.68$ & $2{,}498$ \\
0.25 & 323 & 73  & $+9{,}358$  & $1.97$ & $2{,}064$ \\
0.50 & 272 & 96  & $+10{,}032$ & $2.24$ & $2{,}064$ \\
0.75 & 228 & 82  & $+9{,}305$  & $\mathbf{2.37}$ & $1{,}924$ \\
1.00 & 188 & 82  & $+7{,}776$  & $2.26$ & $1{,}849$ \\
\bottomrule
\end{tabular}
\end{table}

Sharpe peaks at the selected threshold $\theta$ with Sharpe $= 2.37$, net PnL
$= 9{,}305\bp$ over 2024 ($= 93.1\%$ cumulative log return on notional;
$9{,}305 \times 10^{-4} = 0.9305$),
and maximum drawdown $1{,}924\bp$ ($= 19.2\%$). All thresholds are profitable net
of maker costs. The strategy is active on 228 of 365 days (62.5\% exposure)
at the optimal threshold.

\subsection{Direction vs.\ Magnitude}

If N3 is a directional signal (predicting which way BTC will move), then
$\IC(N3, r^{(24h)})$ should exceed $\IC(N3, |r^{(24h)}|)$. If N3 is
a volatility signal (predicting how much BTC will move, but not in which
direction), the reverse holds.

\begin{table}[H]
\centering
\caption{Direction vs.\ magnitude predictability (OOS 2024, daily, 24h)}
\label{tab:dir_mag}
\begin{tabular}{lrr}
\toprule
Target & IC (Pearson) & $|\IC|/\IC^*$ \\
\midrule
Direction: $r^{(24h)}$ & $+0.133$ & $4.87\times$ \\
Magnitude: $|r^{(24h)}|$ & $+0.081$ & $2.03\times$ \\
\bottomrule
\end{tabular}
\end{table}

N3 predicts direction ($4.87\times$) more strongly than magnitude ($2.03\times$).
This confirms N3 is fundamentally a directional signal rather than a
volatility-timing tool, consistent with the fear-resolution mechanism
(Section~\ref{sec:mechanism}).

\subsection{Event Robustness}

The 6 ``extreme event days'' in OOS 2024 are defined as days with
$|r^{(24h)}| > 3\hat{\sigma}_r$ ($> 827\bp$). Examples include the ETF
approval announcement (January 2024) and the halving-period volatility
(April--May 2024).

\begin{table}[H]
\centering
\caption{Event robustness test (OOS 2024, N3, 24h)}
\label{tab:events}
\begin{tabular}{lrrl}
\toprule
Sample & $n$ & IC (Pearson) & $|\IC|/\IC^*$ \\
\midrule
Full sample & 365 & $+0.133$ & $4.87\times$ \\
Ex 6 event days & 359 & $+0.123$ & $4.09\times$ \\
\bottomrule
\end{tabular}
\end{table}

Removing the six largest-return days reduces the IC from $4.87\times$ to
$4.09\times$, a modest change. The signal is not ``made'' by a handful of
outlier events; it is distributed across the full calendar.


%% ============================================================
\section{Kill-Attempt: 2025--2026 Out-of-Sample}
\label{sec:kill}
%% ============================================================

The 2024 OOS results are encouraging, but any single 12-month period
may reflect unusual market conditions. To test whether N3 is a structural
effect or a 2024 artefact, kline and DVOL data were extended through
2026-05-13 and the same non-overlapping daily protocol applied.

\subsection{Walk-Forward IC Across All Periods}

\begin{table}[H]
\centering
\caption{N3 walk-forward IC: all periods (non-overlapping daily, 24h, block-bootstrap $p$).
IC computed via Pearson correlation (Section~\ref{sec:IC}).
The $|\IC|/\IC^*$ ratio uses each period's own $\sigma_r^{(24h)}$.}
\label{tab:n3_wf_all}
\begin{tabular}{lrrrrl}
\toprule
Period & $n$ & IC (Pearson) & $p$ & $|\IC|/\IC^*$ & Verdict \\
\midrule
Train 2023    & 365 & $+0.069$ & $0.204$ & $2.11\times$ & Training \\
OOS 2024-H1   & 182 & $+0.188$ & $0.011$ & $7.05\times$ & \textbf{PASS} \\
OOS 2024-H2   & 184 & $+0.076$ & $0.322$ & $2.70\times$ & Weak \\
OOS 2025-H1   & 181 & $+0.114$ & $0.151$ & $3.72\times$ & Marginal \\
OOS 2025-H2   & 184 & $-0.029$ & $0.684$ & $0.73\times$ & Fail \\
OOS 2026-YTD  & 132 & $+0.024$ & $0.792$ & $0.86\times$ & Fail \\
\bottomrule
\end{tabular}
\end{table}

2024-H1 passes with $p = 0.011$ and ratio $7.05\times$.
2025-H1 is directionally consistent (IC $= +0.114$, ratio $3.72\times$) but
marginal at $p = 0.151$ --- below the 5\% threshold, consistent with the
regime explanation (only 47\% high-DVOL days that period).

2025-H2 and 2026-YTD fail. The signal is intermittent rather than uniformly
persistent.

\subsection{Year-by-Year Backtest}

\begin{table}[H]
\centering
\caption{Year-by-year long-short backtest ($\theta = \theta_\mathrm{sel}$, maker cost, daily)}
\label{tab:yearly}
\begin{tabular}{lrrrrr}
\toprule
Year & $n$ & Sharpe & PnL (bp) & MaxDD (bp) & $n_{\text{trades}}$ \\
\midrule
2023 (train) & 365 & $+0.41$ & $+1{,}650$ & $-1{,}962$ & 77 \\
2024 (OOS)   & 366 & $+1.88$ & $+9{,}375$ & $-2{,}050$ & 84 \\
2025 (OOS)   & 365 & $+0.20$ & $+784$     & $-4{,}638$ & 89 \\
2026 YTD     & 132 & $-0.29$ & $-563$     & $-2{,}230$ & 23 \\
\bottomrule
\end{tabular}
\end{table}

2024 is the dominant year. 2025 is marginally positive but the maximum
drawdown ($-4{,}638\bp$) is much larger than 2024, suggesting the signal
is less consistently reliable. The question is whether this is structural
decay or regime-driven.


%% ============================================================
\section{Regime Conditioning}
\label{sec:regime}
%% ============================================================

\subsection{The IC by DVOL Level}

The following observation is the central diagnostic result of this report.
The OOS 2024+ IC is split by the contemporaneous DVOL level (using the
DVOL observation at signal measurement time, not the $z$-score):

\begin{table}[H]
\centering
\caption{N3 IC by DVOL level regime (OOS 2024+, daily non-overlapping)}
\label{tab:n3_regime}
\begin{tabular}{lrrrr}
\toprule
Regime & $n$ & IC & $|\IC|/\IC^*$ & $p$ \\
\midrule
Low DVOL ($< 51$, below median)   & 430 & $+0.019$ & $0.51\times$ & $0.691$ \\
High DVOL ($\geq 51$, above median) & 433 & $+0.119$ & $4.69\times$ & $0.018$ \\
High DVOL ($\geq 57$, 75th pct) & 216 & $+0.153$ & $6.47\times$ & $0.031$ \\
\bottomrule
\end{tabular}
\end{table}

In low-DVOL environments (DVOL $< 51$, below the median of the full DVOL
history), the IC is statistically and economically indistinguishable from zero
($p = 0.691$, ratio $= 0.51\times$). There is no predictable alpha to capture.

In high-DVOL environments (DVOL $\geq 51$), the signal recovers to $4.69\times$
the maker breakeven ($p = 0.018$). At the 75th percentile (DVOL $\geq 57$),
the IC strengthens further to $6.47\times$ ($p = 0.031$).

This is the empirical confirmation of the fear-resolution mechanism: the
signal only predicts returns when fear is elevated above its local baseline.
In calm markets, there is no fear to resolve, and N3 correctly produces no edge.

\subsection{DVOL Distribution by Period}

This finding directly explains the period-to-period variability in
Table~\ref{tab:n3_wf_all}:

\begin{table}[H]
\centering
\caption{DVOL level distribution by period}
\label{tab:dvol_dist}
\begin{tabular}{lrrrr}
\toprule
Period & $n$ & DVOL mean & DVOL median & Pct $\geq 51$ \\
\midrule
Train 2023   & 365 & 49.5 & 50.5 & 47\% \\
OOS 2024-H1  & 182 & 59.3 & 56.1 & 74\% \\
OOS 2024-H2  & 184 & 56.1 & 56.5 & 88\% \\
OOS 2025-H1  & 181 & 50.5 & 49.8 & 47\% \\
OOS 2025-H2  & 184 & \textbf{41.6} & \textbf{39.5} & \textbf{4\%} \\
OOS 2026-YTD & 133 & 47.5 & 46.1 & 42\% \\
\bottomrule
\end{tabular}
\end{table}

The 2025-H2 failure is explained directly and entirely: the BTC bull market
of mid-2025 suppressed implied volatility to a mean of 41.6 (well below
the 51 regime threshold), with only 4\% of days having DVOL $\geq 51$.
There is no fear premium to trade during a sustained low-volatility
bull market. The signal correctly has nothing to do.

By contrast, 2024-H1 had 74\% of days in the high-DVOL regime and produced
IC $= 7.05\times$ ($p = 0.011$). 2024-H2 had 88\% high-DVOL days; 2025-H1
had 47\% (borderline, Pearson IC $= 3.72\times$, $p = 0.151$).

\textbf{The intermittency is entirely explained by the DVOL level regime.
This is regime dependence, not signal decay.}

\subsection{DVOL Threshold Sweep}
\label{sec:threshold_sweep}

To select the DVOL conditioning threshold without in-sample data snooping,
we evaluate the long-short rule ($\theta = \theta_\mathrm{sel}$, maker cost) over the full
OOS 2024+ period ($n = 863$ days) at varying DVOL filters:

\begin{table}[H]
\centering
\caption{DVOL entry-filter sweep (OOS 2024+, $n=863$ days, $\theta=\theta_\mathrm{sel}$, maker cost)}
\label{tab:dvol_sweep}
\begin{tabular}{rrrrrr}
\toprule
DVOL filter & $n_{\text{active}}$ & IC & $p$ & Sharpe & PnL (bp) \\
\midrule
None       & 863 & $+0.072$ & $0.033$ & $+0.88$ & $+9{,}595$  \\
$\geq 46$  & 590 & $+0.122$ & $0.004$ & $+1.71$ & $+13{,}907$ \\
$\geq 51$  & 445 & $+0.124$ & $0.006$ & $+1.68$ & $+10{,}955$ \\
$\geq 54$  & 333 & $+0.117$ & $0.035$ & $+2.20$ & $+10{,}895$ \\
$\geq 57$  & 210 & $+0.146$ & $0.035$ & $+\mathbf{3.46}$ & $+12{,}246$ \\
$\geq 60$  & 121 & $+0.150$ & $0.120$ & $+3.32$ & $+7{,}472$  \\
\bottomrule
\end{tabular}
\end{table}

Sharpe increases monotonically as the DVOL threshold rises, peaking near
DVOL $\geq 57$ before sample size becomes restrictive (121 active days
at DVOL $\geq 60$ is too few for reliable estimation). The practical sweet
spot lies in the mid-range of the DVOL sweep: Sharpe $2.20$--$3.46$ with
the maximum drawdown shrinking from $5{,}055\bp$ (unfiltered) to $1{,}624\bp$
at the tighter end of the range.

\paragraph{Robustness to the choice of threshold.}
The Sharpe increases smoothly across the sweep range — not concentrated at a
cliff. The selected DVOL threshold $\delta$ is adopted for the
position-level backtest (Section~\ref{sec:backtest}) as the practical
balance between Sharpe improvement and active-day count.

\subsection{Per-Period Conditioned Performance}

Applying the DVOL $\geq 51$ filter by period:

\begin{table}[H]
\centering
\caption{Long-short Sharpe: unconditional vs DVOL $\geq 51$
($\theta=\theta_\mathrm{sel}$, maker cost)}
\label{tab:dvol_cond}
\begin{tabular}{lrrrr}
\toprule
Period & Sharpe (all) & Sharpe ($\geq 51$) & PnL (bp, filtered) & Days passing \\
\midrule
Train 2023   & $+0.41$ & $+1.85$ & $+3{,}812$ & 47\% \\
OOS 2024-H1  & $+2.26$ & $+2.32$ & $+4{,}932$ & 74\% \\
OOS 2024-H2  & $+1.46$ & $+1.51$ & $+3{,}018$ & 88\% \\
OOS 2025-H1  & $+1.10$ & $+1.87$ & $+2{,}187$ & 47\% \\
OOS 2025-H2  & $-1.03$ & $-0.43$ & $-53$       & \textbf{4\%}  \\
OOS 2026-YTD & $-0.29$ & $+0.84$ & $+877$      & 42\% \\
\bottomrule
\end{tabular}
\end{table}

The DVOL filter performs exactly as expected:
\begin{itemize}[noitemsep]
  \item 2025-H2: the filter identifies that only 4\% of days qualify;
        the strategy is essentially flat, correctly avoiding the
        low-signal environment.
  \item 2026-YTD: recovered from Sharpe $-0.29$ to $+0.84$ by
        conditioning on the 42\% of days with DVOL $\geq 51$.
  \item 2025 full year: Sharpe improves from $+0.20$ to $+1.65$; the
        signal is strong in high-IV days even during a structurally
        low-IV year.
\end{itemize}


%% ============================================================
\section{Position-Level Strategy Backtest}
\label{sec:backtest}
%% ============================================================

\subsection{Frozen Rule Specification}

The fully specified trading rule --- \emph{frozen} after Section~\ref{sec:threshold_sweep}
with no further parameter changes --- is:

\begin{enumerate}
  \item \textbf{Entry condition}: N3$_t > \theta$ AND DVOL$_t \geq \delta$.
        Enter LONG. If N3$_t < -\theta$ AND DVOL$_t \geq \delta$, enter SHORT.
        \emph{(Specific threshold values $\theta$ and $\delta$ are withheld.)}
  \item \textbf{Exit}: close position exactly 24 hours after entry.
  \item \textbf{Costs}: maker-side one-way fee of $3\bp$ charged on both
        \emph{entry and exit} of every position.  Every open-and-close cycle
        pays $c_{\text{entry}} + c_{\text{exit}} = 6\bp$ round-trip,
        regardless of whether the close is immediately followed by a new trade.
        No cost is charged on bars where the position is held unchanged.
  \item \textbf{PnL per trade}:
        \[
          \pi_t = \text{pos}_t \times
                  \bigl(r_t^{(24h)} - f_t^{(24h)}\bigr)
                  - c_{\text{entry}} - c_{\text{exit}}
        \]
        where $c_{\text{entry}} = c_{\text{exit}} = 3\bp$ (one-way maker),
        giving a $6\bp$ round-trip deduction per trade.
        In the code: $\texttt{n\_trades} \times \texttt{MAKER\_RT}$
        where $\texttt{MAKER\_RT} = 6\bp$.
  \item \textbf{Rebalancing}: positions checked daily at signal update time.
\end{enumerate}

\subsection{Trade Log and Summary Statistics}

Over 863 OOS calendar days (2024-01-01 to 2026-05-13), the frozen rule
triggered 199 trades:

\begin{center}
\begin{tabular}{lr}
\toprule
Metric & Value \\
\midrule
Trades              & 199 \\
Sharpe (annualised) & $+2.95$ \\
Net PnL             & $+11{,}228\bp$ \\
Max drawdown        & $-1{,}612\bp$ \\
Win rate            & $53.8\%$ \\
Avg winning trade   & $+262\bp$ \\
Avg losing trade    & $-182\bp$ \\
Exposure            & $16.2\%$ of days \\
Profit factor       & $262/182 \times (53.8/46.2) = 1.66$ \\
\bottomrule
\end{tabular}
\end{center}

A Sharpe of $2.95$ with 199 trades over a period that includes both
bull and bear regimes is a strong result. The win rate of $53.8\%$
is modest but the average win ($262\bp$) substantially exceeds the
average loss ($182\bp$), yielding a profit factor of $1.66$.

\subsection{Period Decomposition}

\begin{table}[H]
\centering
\caption{Position-level backtest by period (frozen rule: N3$z > \theta$ AND DVOL $\geq \delta$)}
\label{tab:backtest_periods}
\begin{tabular}{lrrrrrr}
\toprule
Period & $n_{\text{trades}}$ & Sharpe & PnL (bp) & MaxDD (bp) & Win\% & Exp\% \\
\midrule
Train 2023   & 77  & $+2.89$ & $+3{,}641$ & $-1{,}154$ & 51.9\% & 6.3\% \\
OOS 2024-H1  & 81  & $+2.57$ & $+4{,}071$ & $-1{,}496$ & 51.9\% & 6.6\% \\
OOS 2024-H2  & 63  & $+3.01$ & $+3{,}259$ & $-1{,}149$ & 50.8\% & 5.1\% \\
OOS 2025-H1  & 38  & $+2.31$ & $+1{,}769$ & $-1{,}346$ & 65.8\% & 3.1\% \\
OOS 2025-H2  & 4   & $+3.41$ & $+171$     & $-55$      & 50.0\% & 0.3\% \\
OOS 2026-YTD & 13  & $+5.93$ & $+1{,}958$ & $-605$     & 46.2\% & 1.1\% \\
\bottomrule
\end{tabular}
\end{table}

\textbf{Every period shows a positive Sharpe.} This is the strongest result
in this report. The Sharpe is positive in training data, in both 2024 half-years,
in 2025-H1, in 2025-H2 (4 trades only, low-significance), and in 2026 YTD
(13 trades, low-significance). No period shows a negative Sharpe after applying
the DVOL regime filter.

The exposure column is critical to interpretation. The strategy is active
only 6.6\% of days in 2024-H1 and 3.1\% in 2025-H1. The DVOL filter
keeps the strategy flat during the majority of the calendar --- a form
of built-in regime selectivity that substantially reduces the maximum drawdown.

\subsection{Long-Side vs.\ Short-Side Decomposition}
\label{sec:longshort_decomp}

The frozen rule enters both long (N3$_t > +\theta$) and short (N3$_t < -\theta$)
positions conditional on DVOL $\geq \delta$.  The fear-resolution mechanism
(Section~\ref{sec:mechanism}) provides a mechanistically grounded rationale
for longs: an elevated-fear regime resolves toward positive returns as
perceived risk subsides.  The short side (DVOL above $\delta$ but below its
30-day mean) has a weaker theoretical basis and does not directly follow
from the same mechanism.

Table~\ref{tab:longshort_decomp} separates the two sides for OOS 2024
(365 days, 143 signals, maker round-trip cost $= 6\bp$ per trade).

\begin{table}[H]
\centering
\caption{Long vs.\ short performance decomposition (OOS 2024,
  thresh $= \theta$, DVOL $\geq \delta$)}
\label{tab:longshort_decomp}
\begin{tabular}{lrrrrrr}
\toprule
Side & $n$ & Net PnL (bp) & Sharpe & Win\% & Avg win (bp) & MaxDD (bp) \\
\midrule
LONG  ($N3_z > +\theta$) & 105 & $+6{,}080$ & $3.60$ & 52.4\% & $+274$ & $1{,}119$ \\
SHORT ($N3_z < -\theta$) &  38 & $+403$     & $0.81$ & 47.4\% & $+222$ & $1{,}453$ \\
\midrule
Combined               & 143 & $+6{,}483$ & ---    & 51.0\% & ---    & ---       \\
\bottomrule
\end{tabular}
\end{table}

\textbf{The long side accounts for 94\% of net PnL with a Sharpe of $3.60$,
versus $0.81$ for the short side.}
The short side has a win rate below 50\% and a maximum drawdown ($1{,}453\bp$)
that exceeds its total net PnL ($403\bp$) by a factor of $3.6\times$.

\paragraph{Implication for live trading.}
The first live paper-trading version should be \textbf{long-only}: enter
when N3$_t > \theta$ AND DVOL $\geq \delta$, exit after 24 hours.
This retains 94\% of the edge while eliminating the mechanistically weaker
short leg. The short leg should be re-introduced only after accumulating
at least 100 paper-trade observations on the short-only rule.

\subsection{Event Dependence Analysis}

\begin{table}[H]
\centering
\caption{Event robustness: position-level backtest (OOS 2024+)}
\label{tab:backtest_events}
\begin{tabular}{lrrrr}
\toprule
Sample & $n_{\text{trades}}$ & Sharpe & PnL (bp) & MaxDD (bp) \\
\midrule
All OOS days              & 199 & $+2.95$ & $+11{,}228$ & $-1{,}612$ \\
Ex top/bottom 1\% $|r|$  & 194 & $+2.50$ & $+8{,}133$  & $-1{,}659$ \\
Ex 8 known macro events   & 193 & $+2.98$ & $+10{,}767$ & $-1{,}612$ \\
\bottomrule
\end{tabular}
\end{table}

Removing the 8 known macro events (ETF approval, halving, election,
tariff shock, etc.) leaves Sharpe at $2.98$ --- marginally \emph{higher}
than the full sample. The strategy does not depend on these events;
they are approximately neutral to performance. Removing the top and bottom
1\% of return days reduces Sharpe to $2.50$, indicating that some edge
is earned on large-return days, but the strategy comfortably survives
without them.


%% ============================================================
\section{Discussion}
\label{sec:discussion}
%% ============================================================

\subsection{Alternative Mechanisms}

The fear-resolution mechanism (Section~\ref{sec:mechanism}) is the primary
hypothesis. Two alternative explanations deserve consideration:

\paragraph{Options dealer gamma hedging.}
When dealers sell calls (a typical activity in bull markets), they are short
gamma and must buy the underlying as its price rises and sell as it falls,
amplifying intraday moves. When DVOL spikes (consistent with a stress event),
dealers' short-gamma positions cause aggressive hedging, suppressing prices.
When the event passes and dealers unwind hedges, prices recover. This
mechanism is consistent with the positive IC sign but predicts the signal
should be strongest near options expiry cycles (last Friday of each month
for Deribit), a prediction that has not yet been tested.

\paragraph{Perpetual funding crowding.}
When DVOL is elevated, long-side participants (who bought calls or hold
spot) tend to reduce exposure. Perpetual funding rates often fall or
invert during such episodes. When fear resolves, new longs rebuild,
pushing funding positive and exerting upward price pressure. Under this
mechanism, the N3 signal is a proxy for the funding cycle, and the
correlation with DVOL is indirect. The total PnL analysis
(Table~\ref{tab:funding}) cannot distinguish between these mechanisms:
both predict positive IC and are consistent with the funding data.

\subsection{Limitations}

\paragraph{Sample size.}
Each 6-month OOS window contains $n \approx 180$ daily observations.
A block-bootstrap $p$-value of $0.018$ with $B = 1{,}000$ replicates
and $n = 182$ observations is meaningful but not overwhelmingly precise.
The 95\% confidence interval on the IC (roughly $\pm 0.10$) is wide.
The signal needs further confirmation as additional OOS years accumulate.

\paragraph{Transaction cost assumptions.}
Maker-side execution (0.02\% fee, 0.01\% slippage) is assumed throughout.
This requires the strategy to post limit orders that are reliably filled.
For a 24-hour signal with daily rebalancing, limit orders should fill
readily, but in stressed markets (precisely the high-DVOL regime where
the signal is active), execution may slip toward taker costs
($c_{\mathrm{RT}} = 0.10\%$ vs. $0.06\%$). At taker costs, the IC breakeven
rises from $0.028$ to $0.047$ at 24h, reducing the $4.87\times$ ratio to
$\approx 2.8\times$ --- still viable, but with less margin.

\paragraph{Capacity.}
The signal is evaluated on unit-notional positions. BTC perpetual daily
volume on Binance exceeds \$20 billion, so meaningful trading sizes
(up to \$1M notional) would not move the market. At sizes above \$5--10M,
market impact models would need to be applied.

\paragraph{Regime filter stability.}
The selected DVOL $\geq \delta$ threshold is chosen using OOS 2024+ data. While
Table~\ref{tab:dvol_sweep} shows a smooth relationship (not a cliff
at any single value), there is residual data snooping in the threshold choice.
A fixed-rule strategy using DVOL $\geq 51$ (the historical median)
produces Sharpe $1.68$ --- lower but still strongly profitable.

\subsection{Multiple Comparison Adjustment}

This research programme tested 5+ distinct signals across multiple
horizons and parameter configurations. The ``best signal'' N3 is
selected from this set. Under the null hypothesis that no signal is
predictive, the probability that at least one signal passes a 5\%
significance test by chance is substantially larger than 5\%.

Bonferroni correction: if $k = 10$ independent tests are performed,
the per-test threshold for family-wise 5\% error is $\alpha/k = 0.005$.
The N3 24h result at $p = 0.017$ would fail this criterion.

Two arguments partially mitigate this concern:
\begin{enumerate}
  \item N3 passes a 48h test independently ($p = 0.002$), and passes two
        independent 6-month OOS windows. These provide corroborating
        evidence beyond a single $p$-value.
  \item The signals tested are not exchangeable: N1, N2, N3 test
        variations of the same underlying mechanism, not distinct
        hypotheses.
\end{enumerate}

A conservative reader should treat N3 as a strong pilot finding requiring
additional confirmation from future OOS data rather than a conclusively
established edge.


%% ============================================================
\section{Conclusion}
\label{sec:conclusion}
%% ============================================================

We have documented and validated Signal N3 --- the 30-day $z$-score of the
Deribit BTCVOL index --- as a predictor of 24-hour BTCUSDT perpetual futures
returns.

The key results:

\begin{enumerate}
  \item \textbf{Statistical significance}: $\IC = +0.133$, block-bootstrap
        $p = 0.017$ on non-overlapping daily observations (OOS 2024). The
        result is robust to the overlapping-return inflation problem
        (no inflation factor detected). Significance extends to 48h
        ($p = 0.002$).

  \item \textbf{Economic significance}: $|\IC|/\IC^*_{\mathrm{maker}} = 4.87\times$
        at 24h, rising to $8.29\times$ at 48h. Net-of-cost quintile spreads
        of $83.7\bp/\text{day}$ confirm tradability.

  \item \textbf{Regime dependence}: the IC is statistically zero in low-DVOL
        environments (DVOL $< 51$, $p = 0.691$) and strongly positive in
        high-DVOL environments (DVOL $\geq 51$, $p = 0.018$). The
        intermittency observed in 2025-H2 is fully explained by that
        period having only 4\% of days in the high-DVOL regime.

  \item \textbf{Persistence}: 2024-H1 passes the block-bootstrap test
        ($p = 0.011$, ratio $7.05\times$). 2025-H1 is directionally consistent
        (IC $= +0.114$, ratio $3.72\times$, $p = 0.151$) but marginal, consistent
        with its 47\% high-DVOL exposure that period.

  \item \textbf{Position-level performance}: the frozen rule
        (N3$z > \theta$ AND DVOL $\geq \delta$, 24h hold, maker cost) achieves
        Sharpe $+2.95$ over 863 OOS days with maximum drawdown $-1{,}612\bp$
        and positive Sharpe in every sub-period tested.
\end{enumerate}

The primary open question is whether the 2025-H2 and 2026-YTD weakness
represents the beginning of structural decay or is a continuation of the
same regime explanation (insufficient high-DVOL days). Continued OOS data
collection, particularly during a future high-IV regime, will resolve this
question.

\vspace{1em}
\noindent\textit{Data: Binance FAPI 1m klines (2022--2026); Deribit BTCVOL 1h
(2022--2026); Binance 8h funding rates (2022--2026).}

%% ============================================================
\section*{References}
%% ============================================================
\begingroup
\setlength{\parindent}{-2em}
\setlength{\leftskip}{2em}
\setlength{\parskip}{6pt}
\small

Britten-Jones, M., \& Neuberger, A. (2000). Option prices, implied price processes, and stochastic volatility. \textit{The Journal of Finance}, \textit{55}(2), 839--866. \url{https://doi.org/10.1111/0022-1082.00228}

CBOE. (2009). \textit{The CBOE Volatility Index: VIX}. Chicago Board Options Exchange. \url{https://www.cboe.com/micro/vix/vixwhite.pdf}

Jiang, G. J., \& Tian, Y. S. (2005). The model-free implied volatility and its information content. \textit{The Review of Financial Studies}, \textit{18}(4), 1305--1342. \url{https://doi.org/10.1093/rfs/hhi027}

Künsch, H. R. (1989). The jackknife and the bootstrap for general stationary observations. \textit{The Annals of Statistics}, \textit{17}(3), 1217--1241. \url{https://doi.org/10.1214/aos/1176347265}

Liu, R. Y., \& Singh, K. (1992). Moving blocks jackknife and bootstrap capture weak convergence. In R. LePage \& L. Billard (Eds.), \textit{Exploring the Limits of Bootstrap} (pp. 225--248). Wiley.

\endgroup

\end{document}
